% ============================================================
%  ch03_bianchi_bounds.tex — Bianchi Classification and Bounds
%  Requires: macros.tex, amsmath, amssymb, amsthm, booktabs
%  Estimated length: 25–35 pages
%  Rewritten: VD-08, 2026-03-07
% ============================================================

\chapter{Bianchi Type-by-Type Bounds}\label{ch:bianchi-bounds}

The master defect identity of Theorem~\ref{thm:master-defect-id}
provides the algebraic skeleton; the present chapter supplies the
numerical flesh.  We begin with a self-contained account of the
Bianchi classification (Section~\ref{sec:bianchi-class}), then
derive the MES three-bound hierarchy from the covariant Boltzmann
hierarchy (Section~\ref{sec:MES}), specialise the defect identity
to each of the nine types in both orthogonal and tilted sectors
(Sections~\ref{sec:BI}--\ref{sec:classAB}), analyse the Bianchi~V
momentum constraint in detail
(Section~\ref{sec:BV-tilt}), construct the four-dimensional
constraint space (Section~\ref{sec:4D}), and close with the
Frobenius shear--vorticity coupling
(Section~\ref{sec:frobenius-coupling}).


% ============================================================
\section{The Bianchi classification of spatially homogeneous
  spacetimes}\label{sec:bianchi-class}
% ============================================================

\subsection{Spatial Killing vectors and the isometry group}
\label{sec:killing}

Our goal in this section is to classify all spatially
homogeneous---but generically anisotropic---cosmological
geometries.  A spacetime is \emph{spatially homogeneous} if it
admits a group of isometries $G_3$ acting simply transitively
on a family of spacelike hypersurfaces~$\Sigma_t$.  The group
orbits are the surfaces~$\Sigma_t$ themselves: every point on
$\Sigma_t$ can be reached from any other by a group element.  The
isometry group is three-dimensional (matching the dimension of
the spatial sections), and its Lie algebra is specified by the
structure constants $C^c{}_{ab}$ satisfying
\begin{equation}\label{eq:lie-bracket}
[\xi_a, \xi_b] = C^c{}_{ab}\,\xi_c\,,
\end{equation}
where $\{\xi_a\}_{a=1,2,3}$ is a basis for the Killing vector
fields tangent to~$\Sigma_t$.  The antisymmetry
$C^c{}_{ab} = -C^c{}_{ba}$ and the Jacobi identity
$C^e{}_{[ab}C^d{}_{c]e} = 0$ constrain the possible algebras.


\subsection{Decomposition of the structure constants}
\label{sec:struct-decomp}

Ellis \& MacCallum~\cite{EllisMacCallum1969} introduced the
canonical decomposition
\begin{equation}\label{eq:struct-constants}
C^{\alpha}{}_{\beta\gamma}
= \epsilon_{\beta\gamma\delta}\,n^{\alpha\delta}
  + \delta^{\alpha}_{[\beta}\,a_{\gamma]}\,,
\end{equation}
where $n^{\alpha\beta} = n^{(\alpha\beta)}$ is a symmetric
$3 \times 3$ matrix and $a_\alpha$ is a covector.  The Jacobi
identity requires
\begin{equation}\label{eq:jacobi}
n^{\alpha\beta}\,a_\beta = 0\,.
\end{equation}
The symmetric matrix $n^{\alpha\beta}$ can always be
diagonalised; we write its eigenvalues as $(n_1, n_2, n_3)$.
The covector $a_\alpha$ is either zero or, by
Eq.~\eqref{eq:jacobi}, orthogonal to the image of~$n$.

The Bianchi types divide into two classes:
\begin{equation}
\text{Class A:}\quad a_\alpha = 0\,;\qquad
\text{Class B:}\quad a_\alpha \neq 0\,.
\end{equation}
Class~A models have unimodular group actions (the structure
constants satisfy $C^\alpha{}_{\alpha\beta} = 0$); Class~B
models do not.


\subsection{The nine Bianchi types}\label{sec:nine-types}

The possible combinations of $(n_1, n_2, n_3)$ and
$a = |a_\alpha|$ yield nine algebraically distinct Lie algebras.
We collect the complete classification in
Table~\ref{tab:bianchi-types}.

\begin{table}[ht]
\caption{Bianchi classification.  The eigenvalues $(n_1, n_2, n_3)$
of the symmetric part $n^{\alpha\beta}$ and the magnitude
$a = |a_\alpha|$ of the covector determine the type.  The
parameter $h$ in types VI$_h$ and VII$_h$ is the group parameter
$h = a^2/n_1 n_2$ (with $n_3 = 0$); type~III coincides
with~VI$_{-1}$.  The column ``$\Omega_k$ sign'' indicates the sign of
the curvature density parameter $\Omega_k = -\sR/(6\Htheta^2)$ under the
Ellis--MacCallum
convention~\eqref{eq:Rstar-bianchi}: since $\Omega_k = -\sR/(6\Htheta^2)$,
we show the sign of~$\Omega_k$: $0$ (flat), $-$ (open,
$k=-1$ FLRW limit), $+$ (closed, $k=+1$ limit).}
\label{tab:bianchi-types}
\centering\footnotesize\setlength{\tabcolsep}{3.5pt}
\begin{tabular}{llcccccc}
\toprule
Type & Class & $n_1$ & $n_2$ & $n_3$ & $a$
  & $\Omega_k$ sign & FLRW limit \\
\midrule
I      & A & $0$ & $0$ & $0$ & $0$
  & $0$ & $k = 0$ \\
II     & A & $+$ & $0$ & $0$ & $0$
  & $-$ & --- \\
VI$_0$ & A & $+$ & $-$ & $0$ & $0$
  & $-$ & --- \\
VII$_0$& A & $+$ & $+$ & $0$ & $0$
  & $-$ & $k = 0$ (special) \\
VIII   & A & $+$ & $+$ & $-$ & $0$
  & $\pm$ & --- \\
IX     & A & $+$ & $+$ & $+$ & $0$
  & $+$ & $k = +1$ \\
\midrule
III    & B & $+$ & $-$ & $0$ & $\neq 0$
  & $\pm$ & --- \\
IV     & B & $+$ & $0$ & $0$ & $\neq 0$
  & $-$ & --- \\
V      & B & $0$ & $0$ & $0$ & $\neq 0$
  & $-$ & $k = -1$ \\
VI$_h$ & B & $+$ & $-$ & $0$ & $\neq 0$
  & $\pm$ & --- \\
VII$_h$& B & $+$ & $+$ & $0$ & $\neq 0$
  & $-$ & $k = -1$ (special) \\
\bottomrule
\end{tabular}
\end{table}

Three types contain isotropic (FLRW) limits: Bianchi~I
($k = 0$), Bianchi~V ($k = -1$), and Bianchi~IX ($k = +1$),
each reached by setting all anisotropy to zero.  Bianchi~VII$_0$
also admits a flat FLRW limit as a special case, and VII$_h$
admits an open FLRW limit.  These are the types most relevant
for observational cosmology.


\subsection{Kinematic variables by type}\label{sec:kin-by-type}

Not every Bianchi type supports all four kinematic degrees of
freedom.  In the orthogonal sector ($\beta = 0$), the
geometry-frame congruence $u^a$ is the group-orbit normal,
giving $\omega_{ab}^{(G)} = 0$ and $\dot{u}_a^{(G)} = 0$
for all types---the only non-trivial kinematic variable is the
shear $\sigma_{ab}$.  The number of independent shear components
depends on the type (from $1$ for the LRS types to $5$ for the
fully anisotropic types).

Tilting the fluid activates additional variables: the tilt
rapidity $\beta$ and, through the King--Ellis boost
(Section~\ref{sec:king-ellis}), the matter-frame
vorticity $\hat{\omega}_{ab}$ and acceleration
$\hat{\dot{u}}_a$.  The vorticity generated by tilt depends on
the Bianchi type through the structure constants: King \&
Ellis~\cite{KingEllis1973} proved that the tilted fluid has
non-zero vorticity if and only if $n_{\alpha\beta}v^\alpha
\neq 0$, where $v^\alpha$ is the tilt velocity.

For Bianchi~I ($n_{\alpha\beta} = 0$), tilt cannot generate
vorticity regardless of the tilt direction.  For types with
non-degenerate $n_{\alpha\beta}$ (VII$_0$, VIII, IX), vorticity
is generically non-zero for any tilt direction not aligned with
a principal axis.  Type~V ($n_{\alpha\beta} = 0$) produces no
vorticity from the $n$-tensor route, but the non-zero
$a_\alpha$ creates a momentum constraint that algebraically
couples shear to tilt
(Section~\ref{sec:BV-tilt}).

Table~\ref{tab:kin-by-type} summarises the active variables
for each type.

\begin{table}[ht]
\caption{Active kinematic variables by Bianchi type.
``$\sigma$'': shear (present for all types, 1--5 independent
components). ``$\omega$'': matter-frame vorticity generated by
tilt. ``$\dot{u}$'': matter-frame four-acceleration (present
whenever $\beta \neq 0$ and $w \neq 0$). ``$\beta$'': tilt
rapidity. ``C'': tilt constrained by momentum equation.
``F'': tilt free.}
\label{tab:kin-by-type}
\centering\footnotesize\setlength{\tabcolsep}{4pt}
\begin{tabular}{lcccccl}
\toprule
Type & $\sigma$ & $\omega$ (orth.) & $\omega$ (tilted)
  & $\dot{u}$ (tilted) & $\beta$ & Notes \\
\midrule
I      & $\checkmark$ & --- & ---
  & $\checkmark$ & F & No vorticity from tilt \\
II     & $\checkmark$ & --- & $\checkmark$
  & $\checkmark$ & F & 1 tilt component \\
V      & $\checkmark$ & --- & ---
  & $\checkmark$ & C & $\Sigma^2 \propto \beta^2/\Omega_K$ \\
VII$_0$& $\checkmark$ & --- & $\checkmark$
  & $\checkmark$ & F & Generic vorticity \\
VII$_h$& $\checkmark$ & $\checkmark$ & $\checkmark$
  & $\checkmark$ & F & Frobenius: $\Wstd = R^2\Sigstd$ \\
IX     & $\checkmark$ & --- & $\checkmark$
  & $\checkmark$ & F & Generic vorticity \\
\bottomrule
\end{tabular}
\end{table}


% ============================================================
\section{The MES bound hierarchy}\label{sec:MES}
% ============================================================

\subsection{The covariant Boltzmann hierarchy}
\label{sec:boltzmann}

Before diving into the bound derivation, we set up the
Boltzmann equation for the photon distribution function in the
1+3 covariant formalism.  The CMB temperature observed by a
fundamental observer in sky direction $e^a$ (with
$e_a u^a = 0$, $e_a e^a = 1$) is
\begin{equation}\label{eq:T-expansion}
T(x^a, e^b) = \bar{T}(x)\Bigl[
  1 + \sum_{\ell \geq 1}
    \tau_{A_\ell}(x)\,e^{A_\ell}
\Bigr]\,,
\end{equation}
where $\tau_{A_\ell} = \tau_{\langle a_1 a_2 \cdots a_\ell
\rangle}$ are the PSTF temperature anisotropy multipoles and
$e^{A_\ell} = e^{\langle a_1}\cdots e^{a_\ell\rangle}$.  The
dimensionless smallness parameters are
\begin{equation}\label{eq:eps-def}
\eps_\ell := |\tau_{A_\ell}|
= (\tau_{A_\ell}\,\tau^{A_\ell})^{1/2} \ll 1\,.
\end{equation}

The energy-integrated intensity multipoles $I_{A_\ell}$ are
related to the temperature multipoles by $I_{A_\ell}/I \approx
4\,\tau_{A_\ell}$ (at linear order), where
$I \propto \bar{T}^4 \propto \rho_\gamma$ is the radiation
monopole.  The Boltzmann equation, projected into PSTF
multipoles, yields the \emph{multipole divergence equations}
(MDEs)~\cite{MaartensGE1999,ChallinorLasenby1999}:
\begin{equation}\label{eq:MDE-general}
\dot{I}_{A_\ell}
+ \tfrac{4}{3}\Theta\,I_{A_\ell}
+ \tilde{D}^b I_{bA_\ell}
- \frac{\ell}{2\ell+1}\,\tilde{D}_{\langle a_\ell}
  I_{A_{\ell-1}\rangle}
+ \tfrac{4}{3}\,I\,\dot{u}_{a_1}\,\delta_{\ell 1}
- \tfrac{8}{15}\,I\,\sigma_{a_1 a_2}\,\delta_{\ell 2}
= \mathcal{C}_{A_\ell}\,,
\end{equation}
where $\mathcal{C}_{A_\ell}$ denotes Thomson scattering terms.
Two features of this hierarchy are central to the MES programme:

(a) The \textbf{shear} $\sigma_{ab}$ enters \emph{only} in the
$\ell = 2$ equation, through the coefficient $8/15$.

(b) The \textbf{four-acceleration} $\dot{u}_a$ enters
\emph{only} in the $\ell = 1$ equation, through the coefficient
$4/3$.

Vorticity does not appear in the photon energy equation at all
(because $\omega_{ab}e^a e^b = 0$ by antisymmetry); it enters
only through the direction equation and the Einstein constraint
equations, making its bound structurally different from the
other two.


\subsection{Explicit transport equations for $\ell = 0, 1, 2, 3$}
\label{sec:transport}

We write out the first four levels of the hierarchy explicitly
(dropping scattering terms for the free-streaming regime after
decoupling):
\begin{align}
\ell = 0:\quad
& \dot{I} + \tfrac{4}{3}\Theta\,I
  + \tilde{D}^a I_a = 0\,,
  \label{eq:ell0}\\[4pt]
\ell = 1:\quad
& \dot{I}_a + \tfrac{4}{3}\Theta\,I_a
  + \tilde{D}^b I_{ab}
  - \tfrac{1}{3}\tilde{D}_a I
  + \tfrac{4}{3}\,I\,\dot{u}_a = 0\,,
  \label{eq:ell1}\\[4pt]
\ell = 2:\quad
& \dot{I}_{ab} + \tfrac{4}{3}\Theta\,I_{ab}
  + \tilde{D}^c I_{cab}
  - \tfrac{2}{5}\tilde{D}_{\langle a} I_{b\rangle}
  - \tfrac{8}{15}\,I\,\sigma_{ab} = 0\,,
  \label{eq:ell2}\\[4pt]
\ell = 3:\quad
& \dot{I}_{abc} + \tfrac{4}{3}\Theta\,I_{abc}
  + \tilde{D}^d I_{dabc}
  - \tfrac{3}{7}\tilde{D}_{\langle a} I_{bc\rangle}
  = 0\,.
  \label{eq:ell3}
\end{align}
The structure is transparent: each $\ell$ equation couples the
$\ell$-th multipole to its neighbours $\ell \pm 1$ through the
spatial gradient terms $\tilde{D}$.  The kinematic source terms
($\dot{u}_a$ at $\ell = 1$, $\sigma_{ab}$ at $\ell = 2$) break
the otherwise homogeneous hierarchy.


\subsection{The Copernican extension postulate}
\label{sec:copernican}

The MES bounds require not merely that \emph{we} observe
near-isotropy, but that \emph{all} fundamental observers do.
The precise statement, following Stoeger, Maartens \&
Ellis~\cite{StoegerME1995}, is:

\medskip
\noindent\textbf{Copernican Extension Postulate.}\quad
\textit{For all fundamental observers and for all $\ell \geq 1$:
\begin{equation}\label{eq:copernican}
|\tau_{A_\ell}| = O(\eps)\,,\qquad
|\dot{\tau}_{A_\ell}| = O(\eps)\,,\qquad
|\tilde{D}_b\,\tau_{A_\ell}| = O(\eps)\,,
\end{equation}
where $\eps \ll 1$ is the CMB anisotropy parameter.}
\medskip

The derivative-boundedness conditions are essential.  Without
them, the spatial gradient terms in
Eqs.~\eqref{eq:ell1}--\eqref{eq:ell3} are unconstrained, and
the shear inversion from the $\ell = 2$ equation fails.  The
Nilsson--Uggla--Wainwright counter-examples
(Section~\ref{sec:intro-EGS} of Chapter~\ref{ch:intro})
demonstrate precisely this failure mode: almost-isotropic CMB
with unbounded Weyl curvature, violating the derivative
condition.


\subsection{Derivation of the MES shear bound}
\label{sec:MES-shear-derivation}

We now derive the shear bound from the $\ell = 2$
equation~\eqref{eq:ell2}.  Solving for the shear:
\begin{equation}\label{eq:sigma-from-ell2}
\sigma_{ab}
= \frac{15}{8\,I}\Bigl[
  -\dot{I}_{ab} - \tfrac{4}{3}\Theta\,I_{ab}
  - \tilde{D}^c I_{cab}
  + \tfrac{2}{5}\tilde{D}_{\langle a} I_{b\rangle}
\Bigr]\,.
\end{equation}
Under the Copernican postulate, each term on the right-hand
side is $O(\eps)\,\Theta$ (using $I_{ab}/I \sim \eps_2$,
$\dot{I}_{ab}/I \sim \eps_2\,\Theta$, and the gradient
terms involving $\eps_1$ and $\eps_3$).  Taking norms:
\begin{equation}\label{eq:sigma-bound-raw}
\frac{\sqrt{\sigma_{ab}\sigma^{ab}}}{\Theta}
\leq \frac{15}{8}\Bigl[
  c_0\,\eps_2 + c_0'\,\eps_2
  + c_3\,\eps_3 + c_1\,\eps_1
\Bigr]\,,
\end{equation}
where the coefficients $c_0, c_0', c_3, c_1$ arise from the
PSTF Clebsch--Gordan coupling structure of the gradient terms.

In the \textbf{homogeneous case} (Bianchi models), all spatial
gradient terms vanish ($\tilde{D}_a = 0$ on the group orbits),
and the bound tightens dramatically: $\sigma/\Theta \sim
O(\eps_2)$ directly from the quadrupole.  In the
\textbf{inhomogeneous case}, the gradient terms
$\tilde{D}_{\langle a}I_{b\rangle}$ and
$\tilde{D}^c I_{cab}$ contribute terms proportional to
$\eps_1$ and $\eps_3$ respectively, yielding the weaker but
more general bound.

The MES result, combining the numerical prefactors from the
PSTF algebra (the Clebsch--Gordan coefficients for rank-2
coupling to $\ell = 1, 2, 3$), is:

\begin{theorem}[MES shear bound {\normalfont
  $[\ConditionalStatus]$}]\label{thm:MES-shear}
In any almost-FLRW cosmology satisfying:
\begin{enumerate}
\item[\emph{(A1)}] the almost-FLRW linearisation
  ($\sigma/\Theta \ll 1$, $\omega/\Theta \ll 1$),
\item[\emph{(A2)}] the Copernican extension postulate
  (observed multipoles are representative of the full sky),
\item[\emph{(A3)}] the PSTF truncation at $\ell \leq 3$
  (higher multipoles contribute only through the next-order
  coupling, which is suppressed by $\eps_4/\eps_1
  \sim 10^{-3}$),
\end{enumerate}
we have
\begin{equation}\label{eq:Bsig-def}
\frac{\sqrt{\sigma_{ab}\sigma^{ab}}}{\Theta}
< \Bsig
:= \tfrac{5}{3}\,\eps_1 + 3\,\eps_2 + \tfrac{3}{7}\,\eps_3\,.
\end{equation}
\end{theorem}

The coefficients have the following origin:
$5/3$ is the PSTF coupling of the rank-2 shear to the
rank-1 dipole through the gradient
$\tilde{D}_{\langle a}I_{b\rangle}$, involving the
Gaunt coefficient $\langle 1\,0;\,2\,0|1\,0\rangle$;
$3$ is the direct $\ell = 2$ coefficient from the
$I_{ab}$ term in Eq.~\eqref{eq:sigma-from-ell2};
$3/7$ is the coupling to the octupole through
$\tilde{D}^c I_{cab}$, involving the Gaunt coefficient
$\langle 3\,0;\,2\,0|3\,0\rangle$.

\paragraph{Dipole dominance.}
At Planck sensitivity, $\eps_1 \approx 1.234 \times 10^{-3}$
while $\eps_2 \approx 3.56 \times 10^{-6}$ and
$\eps_3 \approx 6.07 \times 10^{-6}$.  The three contributions
to $\Bsig$ are:
\begin{equation}\label{eq:dipole-dominance}
\tfrac{5}{3}\eps_1 = 2.057 \times 10^{-3}\,,\quad
3\eps_2 = 1.068 \times 10^{-5}\,,\quad
\tfrac{3}{7}\eps_3 = 2.600 \times 10^{-6}\,,
\end{equation}
giving $\Bsig = 2.069 \times 10^{-3}$.  The dipole term
contributes $2.057/2.069 = 99.4\%$ of the total.  The shear
bound is, for all practical purposes, a dipole bound.


\subsection{Vorticity and acceleration bounds}
\label{sec:om-udot-bounds}

The vorticity and acceleration bounds follow from the same
Boltzmann hierarchy but through different coupling pathways.

\paragraph{Vorticity bound.}
Vorticity does not appear directly in the photon energy
equation---the contraction $\omega_{ab}e^a e^b$ vanishes
identically by antisymmetry.  The vorticity bound is derived
indirectly, through the Einstein constraint equations.  The
momentum constraint (the $G_{0i}$ field equation) links the
divergence of the shear to the vorticity:
\begin{equation}\label{eq:momentum-constraint}
\tilde{D}^b\sigma_{ab} - \tfrac{2}{3}\tilde{D}_a\Theta
+ \eta_{abc}(\tilde{D}^b\omega^c + 2\dot{u}^b\omega^c)
+ \kappa\,q_a = 0\,.
\end{equation}
Since $\sigma/\Theta = O(\eps)$ and $\dot{u}/\Theta = O(\eps)$
from the shear and acceleration bounds, and the spatial gradients
are bounded by the Copernican postulate, the constraint
propagates the smallness to $\omega/\Theta$.

\begin{theorem}[MES vorticity bound {\normalfont
  $[\ConditionalStatus]$}]\label{thm:MES-omega}
Under the same conditions as Theorem~\ref{thm:MES-shear},
\begin{equation}\label{eq:Bom-def}
\frac{\sqrt{\omega_{ab}\omega^{ab}}}{\Theta}
< \Bom
:= \tfrac{3}{4}\,\eps_1 + 2\,\eps_2 + \tfrac{2}{7}\,\eps_3\,.
\end{equation}
\end{theorem}

The coefficients $(3/4, 2, 2/7)$ arise from the $\Delta\ell =
\pm 1$ coupling of the vorticity vector (rank-1 spatial tensor)
with the photon multipoles, propagated through the constraint
equations.  The vorticity bound is logically \emph{dependent}
on the shear and acceleration bounds---it requires establishing
$\sigma = O(\eps)\Theta$ and $\dot{u} = O(\eps)\Theta$ first,
then using the Einstein constraints to close the system.

\paragraph{Acceleration bound.}
The acceleration enters the $\ell = 1$
equation~\eqref{eq:ell1} through the coefficient $4/3$.
Inverting:
\begin{equation}\label{eq:udot-from-ell1}
\dot{u}_a = \frac{3}{4\,I}\Bigl[
  -\dot{I}_a - \tfrac{4}{3}\Theta\,I_a
  - \tilde{D}^b I_{ab}
  + \tfrac{1}{3}\tilde{D}_a I
\Bigr]\,.
\end{equation}

\begin{theorem}[MES acceleration bound {\normalfont
  $[\ConditionalStatus]$}]\label{thm:MES-udot}
Under the same conditions,
\begin{equation}\label{eq:Bud-def}
\frac{\sqrt{\dot{u}_a\dot{u}^a}}{\Theta}
< \Bud
:= \tfrac{3}{4}\,\eps_1 + \eps_2 + \tfrac{3}{14}\,\eps_3\,.
\end{equation}
\end{theorem}

The coefficient $3/4$ for $\eps_1$ arises from the monopole--dipole
coupling induced by the acceleration source term in
Eq.~\eqref{eq:ell1}; the coefficient $1$ for $\eps_2$ comes from
the gradient term $\tilde{D}^b I_{ab}$; and $3/14$ for $\eps_3$
involves the next-order gradient coupling.


\subsection{Bound ordering hierarchy}\label{sec:hierarchy}

\begin{proposition}[Bound ordering {\normalfont
  $[\ConditionalStatus]$}]\label{prop:ordering}
For any set of non-negative CMB multipoles
$(\eps_1, \eps_2, \eps_3)$:
\begin{enumerate}
\item[\emph{(i)}]
  \textbf{Weak hierarchy} (requires only $\eps_1 > 0$):
  \begin{equation}\label{eq:ordering-weak}
  \Bsig > \Bom \geq \Bud\,.
  \end{equation}
\item[\emph{(ii)}]
  \textbf{Strict hierarchy} (requires $\eps_2 > 0$ or
  $\eps_3 > 0$):
  \begin{equation}\label{eq:ordering-strict}
  \Bsig > \Bom > \Bud\,.
  \end{equation}
\end{enumerate}
\end{proposition}

\begin{proof}
For $\Bsig > \Bom$: the $\eps_1$ coefficients are
$5/3 > 3/4$, so $\eps_1 > 0$ alone suffices.  For
$\Bom \geq \Bud$: the $\eps_1$ coefficients are equal
($3/4 = 3/4$), and the remaining coefficients satisfy
$2 \geq 1$ ($\eps_2$) and $2/7 \geq 3/14$ ($\eps_3$).
This gives part~(i).  Strict inequality $\Bom > \Bud$
requires at least one of $\eps_2 > 0$ or $\eps_3 > 0$,
since the $\eps_2$ coefficient is $2 > 1$ and the $\eps_3$
coefficient is $2/7 > 3/14$.  Planck measures $\eps_2 > 0$,
so part~(ii) holds observationally.
\end{proof}

\begin{remark}[Near-degeneracy of the lower hierarchy]
\label{rem:lower-degeneracy}
The gap between the two lower bounds is
$\Bom - \Bud = \eps_2 + \eps_3/14$.
Strictness therefore requires only that $\eps_2 > 0$
\emph{or} $\eps_3 > 0$; the condition $\eps_2 > 0$ stated
above is sufficient but not necessary.  In the observational
regime $(\eps_2, \eps_3) \ll \eps_1$, this gap is
$\sim 4 \times 10^{-6}$, so the lower hierarchy
$\Bom \gtrsim \Bud$ is numerically near-degenerate even
when strictly ordered.
\end{remark}

The physical content is that the CMB is most
sensitive to shear (rank-2, direct quadrupole source) and least
sensitive to acceleration (rank-1, entering through
monopole--dipole coupling).  The hierarchy is preserved under
dynamical evolution
(Remark~\ref{rem:hierarchy-conservation} of
Chapter~\ref{ch:framework}).
the conditioned legacy figure gallery in Appendix~\ref{app:conditioned-legacy-figure-gallery} shows the three bounds as
functions of the observed multipole amplitudes $(\eps_1,
\eps_2, \eps_3)$, with the strict hierarchy
$\Bsig > \Bom > \Bud$ visible across the full observational range.


\subsection{Standardised-variable form and frame correction}
\label{sec:std-form}

From the definitions $\Sigstd = \sigma_{ab}\sigma^{ab}/(6\Htheta^2)$
and $\Htheta = \Theta/3$, we have
$(\sigma/\Theta)^2 = (2/3)\Sigstd$, giving:
\begin{equation}\label{eq:Sig2-from-Bsig}
\Sigstd < \tfrac{3}{2}\,\Bsig^2\,,\qquad
\Wstd < \tfrac{3}{2}\,\Bom^2\,,\qquad
\Astd < \tfrac{3}{2}\,\Bud^2\,.
\end{equation}

For tilted models, the observed $\eps_1$ contains the
acceleration contribution
$\eps_1^{(\dot{u})} = \etaud\beta$
(Section~\ref{sec:accel-dipole}).  Using the corrected
safe-route $\beta \leq \eps_1/(1 + \etaud)$ and the fact that
$\Bsig$ is linear in $\eps_1$, the frame-corrected bound is
\begin{equation}\label{eq:Bsig-corr}
\Bsig^{\mathrm{corr}}
= \Bsig(\eps_1) + \tfrac{5}{3}\,\etaud\,\beta
\approx (1 + 2.69\,\eps_1)\,\Bsig\,.
\end{equation}
The correction factor $R = 1 + 2.69\eps_1$ produces a
$\sim 0.4\%$ effect on $\Bsig$ and a $\sim 0.8\%$ effect on
$\Sigstd$, propagating to $\sim 16\%$ on $\xmax$ for
shear-dominated types (since $\xmax \propto \Bsig^2$).

The frame-attribution bias $\Bbias = \Bsig^{\mathrm{corr}} -
\Bsig$ must be reported alongside $\Bsig$ values in all
subsequent analysis, per the VT-07 mandate.


% ============================================================
\section{Bianchi~I: the minimal benchmark}\label{sec:BI}
% ============================================================

Bianchi~I is the simplest anisotropic cosmology: flat spatial
sections with no vorticity and no spatial curvature.  The
structure constants vanish entirely ($n^{\alpha\beta} = 0$,
$a_\alpha = 0$), giving $\sR = 0$ and $\Omega_k = 0$.

\begin{theorem}[Bianchi~I defect identity {\normalfont $[\ExactStatus]$}]
\label{thm:BI-defect}
For orthogonal Bianchi~I ($\beta = 0$), $\xI = \Sigstd$.
For tilted Bianchi~I, $\xI = \Sigstd + \Otilt$.
\end{theorem}

\begin{proof}
We apply Theorem~\ref{thm:master-defect-id} with the flat
comparator ($\Okref = 0$).  Since $\Omega_k = 0$ for type~I,
$\Okaniso = \Omega_k - \Okref = 0$.  In the orthogonal sector,
$\beta = 0$ gives $\Otilt = 0$, and $\omega_{ab}^{(G)} = 0$
gives $\Wstd = 0$.  The identity reduces to
$\xI = \Sigstd$.  For tilted type~I, $\Otilt =
(1+w)\Omega_m\sinh^2\!\beta > 0$ while $\Wstd$ remains zero
(type~I cannot generate matter-frame vorticity because
$n_{\alpha\beta} = 0$).  Hence $\xI = \Sigstd + \Otilt$.
\end{proof}

The observational bound follows immediately:
$\xI < (3/2)\Bsig^2$ (orthogonal) or
$\xI < (3/2)[\Bsig^{\mathrm{corr}}]^2 + \Otilt(\betasr)$
(tilted).  At the S3 scenario, the orthogonal bound gives
$\xI < 6.42 \times 10^{-6}$.

Bianchi~I serves as the \emph{minimal anisotropy benchmark}
in the irrotational non-negative sector: for flat orthogonal
types (BVII$_0$, BII, BVI$_0$) under the flat comparator,
$x = \Sigstd = \xI$.  This ordering does \emph{not} hold
generally: types with negative $\Okaniso$ (BVIII, BIX under
the flat comparator) can have $x < 0 < \xI$, and tilted BI
itself has $x = \Sigstd + \Otilt > \xI$ (the tilt raises the
defect above the pure-shear level).

For tilted Bianchi~I, the geometry-frame vorticity vanishes
($\omega_{ab}^{(G)} = 0$) because the structure constants are
trivial.  The matter-frame vorticity is also zero in LRS
(locally rotationally symmetric) tilted type~I, but can be
non-zero in the generic (non-LRS) tilted case when the tilt
axis is not aligned with a shear eigenaxis~\cite{King1973}.
We restrict to the LRS case throughout this analysis.


% ============================================================
\section{Tilted Bianchi~V and the momentum constraint}
\label{sec:BV-tilt}
% ============================================================

Bianchi~V occupies a distinguished position: it is the unique
type for which tilt is not a free parameter but is
\emph{geometrically constrained} by the momentum equations.
This constraint leads to a quadrupole overproduction that
excludes Bianchi~V as a viable cosmological model at any tilt
consistent with the observed dipole anomaly.


\subsection{Geometric structure}\label{sec:BV-geom}

The type~V structure constants are $n^{\alpha\beta} = 0$ and
$a_\alpha \neq 0$ (Class~B).  The spatial Ricci scalar is
$\sR = 4\,a_\alpha a^\alpha > 0$, corresponding to open
spatial sections with $\Omega_k = -\sR/(6\Htheta^2) < 0$.
The isotropic limit ($\sigma_{ab} = 0$, $\beta = 0$) is the
$k = -1$ FLRW model.

A key algebraic property distinguishes type~V from all other
curved types: the traceless spatial Ricci tensor ${}^*\!S_{ab}$
vanishes identically:
\begin{equation}\label{eq:BV-Sab}
{}^*\!S_{ab} := {}^*\!R_{\langle ab\rangle}
= {}^*\!R_{ab} - \tfrac{1}{3}\sR\,h_{ab} = 0\,.
\end{equation}
This vanishing follows because $n^{\alpha\beta} = 0$: the spatial
curvature is isotropic.  The consequence for the defect identity
is immediate: using the open comparator
$\Okref = \Omega_k^{(\mathrm{BV})}$, we have $\Okaniso = 0$.
The defect reduces to
$\xV = \Sigstd - \Wstd + \Otilt$.
Since type~V with $n_{\alpha\beta} = 0$ generates no
matter-frame vorticity from tilt (the King--Ellis condition
$n_{\alpha\beta}v^\alpha = 0$ is trivially satisfied), we have
$\Wstd = 0$ and
\begin{equation}\label{eq:xV-simple}
\xV = \Sigstd + \Otilt\,.
\end{equation}


\subsection{The momentum constraint}\label{sec:BV-momentum}

The Einstein $G_{0i}$ field equations for Bianchi~V impose a
constraint that algebraically couples the shear to the tilt.
In the DCP notation of
Krishnan~et~al.~\cite{Krishnan2022,Krishnan2023}, the
momentum constraint reads
\begin{equation}\label{eq:BV-momentum-DCP}
\frac{2A_0}{X}\Bigl(\frac{\dot{X}}{X}
  - \frac{\dot{Y}}{Y}\Bigr)
= (\rho + p)\,\sinh\beta\,\cosh\beta\,.
\end{equation}
The left-hand side is proportional to the shear
($\sigma_{\mathrm{DCP}} = \dot{X}/X - \dot{Y}/Y$), and
$A_0$ encodes the spatial curvature ($\Omega_K \propto
A_0^2/X^2$).  We now convert to standardised variables.

The LRS Bianchi~V shear satisfies
$\sigma_{ab}\sigma^{ab} = 2\sigma^2 = \frac{2}{3}\sigma_{\mathrm{DCP}}^2$
(the factor $2/3$ from the diagonal traceless shear in three
dimensions), so $\Sigstd = \sigma_{\mathrm{DCP}}^2/(6\Htheta^2)$.
From Eq.~\eqref{eq:BV-momentum-DCP}, squaring both sides and
converting:
\begin{equation}\label{eq:BV-Sig2-exact}
\Sigstd
= \frac{[(1+w)\,\Omega_m]^2\,\sinh^2\!\beta\,\cosh^2\!\beta}
       {4\,|\Omega_K|}\,.
\end{equation}
For $\beta \ll 1$ (valid at the dipole-anomaly amplitude
$\beta \sim 10^{-3}$), $\sinh\beta\,\cosh\beta \approx \beta$,
and the constraint simplifies to
\begin{equation}\label{eq:BV-momentum}
\boxed{
\Sigstd
= \frac{[(1+w)\,\Omega_m]^2\,\beta^2}{4\,|\Omega_K|}
}\,.
\end{equation}
Note that $\Sigstd \propto \beta^2/|\Omega_K|$: the shear is
quadratic in the tilt and inversely proportional to the spatial
curvature.  For small curvature ($|\Omega_K| \ll 1$), a modest
tilt generates enormous shear.


\subsection{Quadrupole overproduction}\label{sec:BV-quad}

We now demonstrate that the momentum
constraint~\eqref{eq:BV-momentum} forces Bianchi~V to
overproduce the CMB quadrupole by a catastrophic margin.

At the CF4 tilt rapidity $\beta = 1.334 \times 10^{-3}$
(WFH2009 primary value; the CF4++ sensitivity variant
$\beta = 1.051 \times 10^{-3}$ gives
$\Sigstd \approx 2.74 \times 10^{-5}$---qualitatively
identical)
and the Planck-preferred curvature
$|\Omega_K| = 0.001$, with $w = 0$ (dust) and
$\Omega_m = 0.3153$:
\begin{equation}\label{eq:BV-Sig2-num}
\Sigstd
= \frac{(0.3153)^2 \times (1.334 \times 10^{-3})^2}
       {4 \times 0.001}
= \frac{1.773 \times 10^{-7}}{4 \times 10^{-3}}
\approx 4.43 \times 10^{-5}\,.
\end{equation}
The MES algebraic ceiling is
$\Sigmax = (3/2)\Bsig^2 \approx 6.42 \times
10^{-6}$.  The momentum-constrained value exceeds the MES
ceiling by a factor of $\sim 7$.

The situation is far worse when we compute the implied CMB
quadrupole.  Using the vector-mode transfer function
$T_2 = 5{,}247.5\;\mu\mathrm{K}/(\sigma/H)$ (VER05 corrected; see Appendix), the shear-induced quadrupole is
$D_2^{\mathrm{shear}} \sim T_2^2 \times \Sigstd$,\footnote{The pre-VER05 formula incorrectly included a factor of $T_0$; the corrected formula is $D_2^{\mathrm{shear}} = (6/2\pi)(T_2\,\sigma_H)^2$ where $T_2$ is in $\mu\mathrm{K}/(\sigma/H)$.}\, where $T_0^2
\sim 10^{15}\;\mu\mathrm{K}^2$---a $10^{13}$-fold
overproduction relative to the observed
$D_2^{\mathrm{obs}} = 225.9\;\mu\mathrm{K}^2$.

This overproduction is independent of the equation of state:
for radiation ($w = 1/3$), the prefactor $(1+w)^2$ increases by
$78\%$, worsening the exclusion.  Bianchi~V is
\emph{falsified by geometry}: there is no parameter-space region
that simultaneously accommodates the CF4 bulk flow and the
observed quadrupole.

\begin{remark}[Tilt bound from the momentum constraint
  {\normalfont $[\ConditionalStatus]$}]
\label{rem:BV-tilt-bound}
Under the following conditions---LRS (locally rotationally
symmetric) geometry, single perfect fluid ($w = 0$), DCP
(deparametrised Copernican principle) dictionary connecting
$\Sigstd$ to $D_2$, and small-$\beta$ linearisation---the
momentum constraint~\eqref{eq:BV-momentum} can be read as
an \emph{upper bound on $\beta$} for a given $\Sigstd$:
$\beta \leq 2\sqrt{|\Omega_K|\,\Sigstd}/[(1+w)\Omega_m]$.
At the MES ceiling $\Sigstd = (3/2)\Bsig^2$ and the Planck-preferred
$|\Omega_K| = 0.001$, the maximum tilt is
$\beta_{\mathrm{BV}} \approx 5.1 \times 10^{-4}$,
which is $\sim 2.2$ times tighter than the safe-route bound
$\betasr = 1.14 \times 10^{-3}$.  The constraint becomes
dramatically stronger for smaller $|\Omega_K|$: at
$|\Omega_K| = 10^{-4}$, the ratio rises to $\sim 7$.  In all
cases, the momentum constraint effectively removes Bianchi~V
from the set of viable tilted cosmologies because the
\emph{implied shear} at the CF4 tilt rapidity exceeds the MES
ceiling by a factor of $\sim 7$.
\end{remark}


% ============================================================
\section{Bianchi~VII$_h$: the richest geometry}
\label{sec:BVIIh}
% ============================================================

Type~VII$_h$ admits shear, vorticity, curvature, and a
``spiral'' parameter $x_h = \sqrt{h}/\sqrt{|\Omega_K|}$
that controls the scale of the helical pattern in the CMB.
The structure constants have $n^{\alpha\beta} =
\mathrm{diag}(n_1, n_2, 0)$ with $n_1, n_2 > 0$, and
$a_\alpha \neq 0$ with $h = a^2/(n_1 n_2)$ defining the
group parameter.


\subsection{Vorticity--shear coupling (Frobenius)}
\label{sec:VIIh-frobenius}

The Jacobi identities for type~VII$_h$ impose an algebraic
relation between the vorticity and shear.  For the LRS
(locally rotationally symmetric) sector, the Wainwright--Ellis
constraint gives
\begin{equation}\label{eq:WS-coupling}
\Wstd = R_{\mathrm{WS}}^2\,\Sigstd\,,
\end{equation}
where $R_{\mathrm{WS}}$ depends on the ratio $h/|\Omega_K|$.
For the observationally preferred range
($h/|\Omega_K| \in [0.1, 10]$),
$R_{\mathrm{WS}} \in [0.9, 1.2]$; the reference value used
throughout this work is $R_{\mathrm{WS}} = 1.06$.

The coupling~\eqref{eq:WS-coupling} means that any constraint
on the vorticity automatically tightens the shear bound and
vice versa.  The Saadeh 95\% upper limit
$(\omega/H)_0 < 5.2 \times 10^{-11}$, translating to
$\Wstd < 9.01 \times 10^{-22}$
(Eq.~\eqref{eq:omH-W2}), then implies
\begin{equation}\label{eq:Sig-from-Saa}
\Sigstd < \frac{\WSaa}{R_{\mathrm{WS}}^2}
= \frac{9.01 \times 10^{-22}}{1.12}
\approx 8.04 \times 10^{-22}\,.
\end{equation}
Equation~\eqref{eq:Sig-from-Saa} is the tightest model-dependent constraint on shear in
the entire Bianchi classification, achieved not by a direct
shear measurement but by the geometric coupling that transmits
the vorticity bound to the shear sector---six orders of
magnitude tighter than the MES algebraic bound.


\subsection{Decay and growing modes}\label{sec:VIIh-modes}

The regular tensor mode in VII$_h$ admits two behaviours: a
decaying mode with transfer function
$T_2^{\mathrm{decay}} = 5{,}247.5\;\mu\mathrm{K}/(\sigma/H)$ and a growing
mode with $T_2^{\mathrm{grow}} = 1.05\;\mu\mathrm{K}/(\sigma/H)$
(VER05 BASS-calibrated values; $\pm 20\%$ calibration
uncertainty)~\cite{Saadeh2016b,Pontzen2011}.

The decaying mode produces a detectable CMB signal at
$(\sigma/H)_0 \sim 10^{-10}$, well below the MES algebraic
ceiling.  The growing mode, by contrast, requires
$(\sigma/H)_0 \sim 10^{-6}$ for a comparable signal---close
to the MES bound itself.  This growing mode is the unique
solution compatible with inflationary initial conditions
(Collins \& Hawking~\cite{CollinsHawking1973}), but its
associated vorticity (through the Frobenius coupling) exceeds
the Saadeh bound by orders of magnitude, leading to decisive
exclusion ($\ln\mathcal{B} = -18.7$) in the Bayesian analysis
(Chapter~\ref{ch:evidence}).

The Saadeh--MES complementarity is most visible for this type.
The $\sim 6$ order-of-magnitude gap between the Saadeh bound
($\Sigstd < 8 \times 10^{-22}$) and the MES algebraic bound
($\Sigstd < 6.4 \times 10^{-6}$) shrinks to approximately one
order for the regular tensor growing mode, where the transfer
function is $\sim 5000\times$ weaker and the dynamical lever arm
collapse.  The two frameworks are complementary: MES provides
type-independent algebraic ceilings; Saadeh~et~al.\ exploit
the full morphological information of the Planck sky.


\subsection{Tilted VII$_h$}\label{sec:VIIh-tilted}

Adding tilt to VII$_h$ introduces a four-parameter model
$(\Sigstd, \beta, \Wstd, x_h)$ that is the most complex in
the classification.  The vorticity Occam cost
($\sim 1.4$~nats) reduces the evidence relative to simpler
tilted types, but the model remains preferred over FLRW only in
the legacy transfer-conditional ranking
($\ln\mathcal{B} \approx +24$) in the evidence analysis
of Chapter~\ref{ch:evidence}.


% ============================================================
\section{Bianchi~IX: the closed model}\label{sec:BIX}
% ============================================================

Type~IX has closed spatial sections ($\sR < 0$ in the
Ellis--MacCallum convention, corresponding to $\Omega_k > 0$)
and admits a natural closed comparator
$\Okref = \Omega_k^{(\mathrm{IX})}$.


\subsection{Curvature split}

For the curvature-matched comparator, $\Okaniso = 0$ and the
defect reduces to $\xIX = \Sigstd + \Otilt$, identical in
form to Bianchi~I.  The distinction lies in the \emph{indirect}
effect of curvature: the Planck constraint
$|\Omega_K| < 0.002$ (95\% CL) limits the curvature contribution
to $\Okaniso \lesssim 10^{-3}$ under the flat comparator, but
under the matched comparator $\Okaniso$ vanishes by definition.

The Collins--Hawking theorem guarantees that all type~IX models
with $\Lambda > 0$ isotropise at late times.  In the concordance
model, the present-epoch shear is exponentially suppressed below
the MES bound, and the observational constraint is driven by the
tilt rapidity through the dipole anomaly.


\subsection{The orthogonal bound}

For orthogonal type~IX (no tilt), the defect with the flat
comparator is $\xIX = \Sigstd + \Okaniso$, where
$\Okaniso = \Omega_k > 0$.  Using the MES shear ceiling and the
Planck curvature prior:
\begin{equation}
\xIX < (3/2)\Bsig^2 + |\Omega_K|
\approx 6.42 \times 10^{-6} + 2 \times 10^{-3}
\approx 2 \times 10^{-3}\,.
\end{equation}
The curvature term dominates in the flat-comparator analysis.
Under the matched comparator, the bound drops to the
Bianchi~I value.


% ============================================================
\section{Remaining types}\label{sec:classAB}
% ============================================================

The six types not yet discussed decompose into two groups.

\subsection{Class~A types (II, VI$_0$, VII$_0$, VIII)}

Types~II, VI$_0$, VII$_0$, and~VIII are Class~A
($a_\alpha = 0$).  Each admits a curvature-matched comparator
for which $\Okaniso = 0$.  Their orthogonal defect bounds are
therefore identical to Bianchi~I up to the curvature matching.
For the flat comparator, types with non-zero $n^{\alpha\beta}$
acquire an $\Okaniso$ term that is generically small
($\lesssim 10^{-3}$) under Planck curvature priors.

In the tilted sector, all four types admit free tilt ($\beta$ is
not constrained by a momentum equation).  Types~II, VII$_0$,
and~VIII generate matter-frame vorticity through the tilt
(since $n_{\alpha\beta} \neq 0$), while type~VI$_0$ (with
$n_1 n_2 < 0$) admits vorticity generation for specific tilt
directions.


\subsection{Class~B types (III, IV, VI$_h$)}

Types~III, IV, and~VI$_h$ are Class~B ($a_\alpha \neq 0$).
The non-vanishing covector introduces a momentum constraint
that couples shear to tilt.  For type~V, this coupling is
algebraically exact (Section~\ref{sec:BV-tilt}).  For type~III
($= \mathrm{VI}_{-1}$), the constraint is weaker, allowing
independent variation of $\Sigstd$ and $\beta$.  Type~IV is
cosmologically marginal: it does not admit an FLRW limit and
is not considered further.


% ============================================================
\section{Summary of type-by-type bounds}\label{sec:summary}
% ============================================================

Table~\ref{tab:type-summary} collects the defect bounds for
all types under the S3 observational scenario, distinguishing
orthogonal and tilted sectors.

\begin{table}[ht]
\caption{Defect bounds by Bianchi type (S3 scenario,
$\eps_1 = 1.476 \times 10^{-3}$).  ``Orth.\ $x_{\max}$'':
orthogonal bound from MES algebraic ceiling.  ``Tilt'': F~=~free
parameter, C~=~constrained by momentum equation.
``Frobenius'': whether the Wainwright--Ellis
coupling~\eqref{eq:WS-coupling} applies.
All bounds include the VT-07 frame correction.}
\label{tab:type-summary}
\centering\footnotesize\setlength{\tabcolsep}{3.5pt}
\begin{tabular}{llccll}
\toprule
Type & Class & $\Omega_k$ sign
     & Orth.\ $x_{\max}$ & Tilt & Frobenius \\
\midrule
I       & A & $0$  & $(3/2)\Bsig^2$         & F  & --- \\
II      & A & $-$  & $(3/2)\Bsig^2$         & F  & --- \\
VI$_0$  & A & $-$  & $(3/2)\Bsig^2$         & F  & --- \\
VII$_0$ & A & $-$  & $(3/2)\Bsig^2$         & F  & --- \\
VIII    & A & $\pm$& $(3/2)\Bsig^2 + \Okaniso$ & F  & --- \\
IX      & A & $+$  & $(3/2)\Bsig^2 + \Okaniso$ & F  & --- \\
\midrule
III     & B & $\pm$& $(3/2)\Bsig^2 + \Okaniso$ & F  & --- \\
V       & B & $-$  & $(3/2)\Bsig^2$             & C  & --- \\
VII$_h$ & B & $-$  & $\WSaa/R_{\rm WS}^2$       & F  &
  $\Wstd = R_{\rm WS}^2\Sigstd$ \\
\bottomrule
\end{tabular}
\end{table}

The table reveals a clear trichotomy.  Types without curvature
(I, VII$_0$) have the simplest bounds.  Types with curvature
(II, III, V, VI$_0$, VIII, IX) acquire an $\Okaniso$ term
that is small under Planck priors.  Type~VII$_h$ alone has a
Frobenius coupling that tightens the bound by six orders of
magnitude.  the conditioned legacy figure gallery in Appendix~\ref{app:conditioned-legacy-figure-gallery} presents the
per-type evidence and departure summary that results from
applying this bound hierarchy to data.

\paragraph{Robustness.}

The type-by-type bounds in this chapter depend on the transport
hierarchy only through the orthogonal Thomson anchor and its
tilted-frame generalisation.  The latter is now isolated in the
Layer-B construction of Section~\ref{sec:tilted-thomson-layer-b},
where the $\beta \to 0$ byte-identity reduction to the LB-4
collision operator is enforced explicitly and the first
polarisation-rotation corrections are quantified on the same
PSTF basis used here.


% ============================================================
\section{The four-dimensional constraint space}
\label{sec:4D}
% ============================================================

\subsection{Constraint inventory}\label{sec:constraint-inventory}

The full constraint space is the region in
$(\Sigstd, \Wstd, \Astd, \beta)$ compatible with all
observational bounds.  We collect the constraints:
\begin{equation}\label{eq:constraint-inventory}
\begin{aligned}
\text{MES shear:} &\quad \Sigstd < \tfrac{3}{2}\Bsig^2\,,\\
\text{MES vorticity:} &\quad \Wstd < \tfrac{3}{2}\Bom^2\,,\\
\text{MES acceleration:} &\quad \Astd < \tfrac{3}{2}\Bud^2\,,\\
\text{Saadeh vorticity:} &\quad \Wstd < 9.01 \times 10^{-22}\,,\\
\text{Safe-route tilt:} &\quad \beta < \betasr\,,\\
\text{Frobenius (VII$_h$):} &\quad
  \Wstd = R_{\mathrm{WS}}^2\,\Sigstd\,,\\
\text{Momentum (V):} &\quad
  \Sigstd = [(1{+}w)\Omega_m]^2\beta^2/(4|\Omega_K|)\,.
\end{aligned}
\end{equation}


\subsection{Six-panel atlas}\label{sec:atlas}

The six independent two-dimensional projections of the
four-dimensional space are:
$(\Sigstd, \Wstd)$, $(\Sigstd, \Astd)$, $(\Sigstd, \beta)$,
$(\Wstd, \Astd)$, $(\Wstd, \beta)$, and $(\Astd, \beta)$.
the conditioned legacy figure gallery in Appendix~\ref{app:conditioned-legacy-figure-gallery} presents this atlas at the
S3 scenario.

The most striking feature is the factor $\sim 10^6$ between the
MES algebraic bound on $\Wstd$ and the Saadeh model-dependent
bound, visible in the $(\Sigstd, \Wstd)$ panel as a narrow
horizontal strip within the larger box.  The Frobenius coupling
for type~VII$_h$ appears as a diagonal line
$\Wstd = R_{\mathrm{WS}}^2\Sigstd$ that intersects the Saadeh
strip at $\Sigstd \approx 8 \times 10^{-22}$---the
model-dependent shear ceiling for this type.

The $(\Sigstd, \beta)$ panel shows the Bianchi~V momentum
constraint as a parabola $\Sigstd \propto \beta^2$, lying
entirely above the MES shear ceiling for $\beta > \beta_{\mathrm{BV}}
\approx 7.3 \times 10^{-5}$.  All other tilted types occupy the
rectangular region bounded by $\Sigstd < (3/2)\Bsig^2$ and
$\beta < \betasr$.


% ============================================================
\section{Frobenius coupling and tilt-generated vorticity}
\label{sec:frobenius-coupling}
% ============================================================

\subsection{The shear--vorticity relation}
\label{sec:shear-vort-relation}

For Bianchi types admitting both shear and vorticity (types
VII$_h$ and, in the tilted sector, types V, III, and others),
the Jacobi identities impose algebraic relations between the
kinematic variables.  The relation~\eqref{eq:WS-coupling} for
type~VII$_h$ is the strongest, but similar couplings arise for
other types through the structure constants.

The physical origin of this coupling is the requirement that
the spatial curvature be consistent with the Bianchi Lie algebra.
The vorticity and shear are not independent geometric degrees of
freedom in a spatially homogeneous spacetime: they are both
constrained by the structure constants through the constraint
equations.  The Frobenius coupling makes this constraint explicit
in standardised variables.


\subsection{Tilt-generated vorticity}\label{sec:tilt-vort}

In tilted Bianchi models, the King--Ellis boost generates
matter-frame vorticity even in types that are irrotational in
the orthogonal sector (Section~\ref{sec:king-ellis}).  The
mechanism proceeds through the antisymmetric part of the
boosted velocity gradient: the leading-order result is
\begin{equation}\label{eq:tilt-vort}
\hat{\omega}_{ab} \propto \gamma^2\,
  v_{[a}\,\sigma_{b]c}\,v^c\,,
\end{equation}
which vanishes when $\sigma_{ab} = 0$ (no shear) or when $v^a$
is a shear eigenvector.  For type~I, this vanishes trivially
(there is no structure-constant contribution to redirect the
tilt), but for types~II, VII$_0$, VIII, and~IX, the
tilt-generated vorticity is generically non-zero.

The magnitude of this secondary vorticity scales as
$|\hat{\omega}| \sim \beta\,\sigma/\Theta$ in the linearised
regime.  At the dipole-anomaly amplitudes
($\beta \sim 10^{-3}$, $\sigma/\Theta \sim 10^{-3}$),
the generated vorticity is $|\hat{\omega}|/\Theta \sim 10^{-6}$,
which exceeds the Saadeh bound by four orders of magnitude
\emph{but is below the MES algebraic bound}.  Whether this
vorticity is detectable depends on the specific Bianchi type and
the morphological signature it produces in the CMB, an analysis
pursued by the Saadeh~et~al.\ template-fitting programme.


% ============================================================
\section{Certification classes and the obstruction theorem}
\label{sec:certification}
% ============================================================

The results of the preceding sections show that the
filling fraction $\FC = x/\xmaxC$ is not universally
well-defined across the full Bianchi classification.  We now
formalise this into an obstruction theorem.

\begin{definition}[Certification class]
\label{def:cert-class}
A Bianchi model belongs to certification class
$\mathcal{G} \in \{\text{certifiable}, \text{restricted},
\text{blocked}\}$ according to the sign-domain sector of
Proposition~\ref{prop:x-sign}:
\begin{enumerate}
\item \emph{Certifiable}: irrotational non-negative sector
  ($\Wstd = 0$, $x \geq 0$).  The standard $\FC \in [0,1]$
  is well-defined.
\item \emph{Restricted}: irrotational negative sector
  ($\Wstd = 0$, $x < 0$).  A signed saturation coordinate
  $\FC^{\pm} = x/|\xmaxC|$ is required.
\item \emph{Blocked}: vortical sector ($\Wstd > 0$).
  No certification route exists without additional
  proof that $\Sigstd - \Wstd + \Otilt \geq 0$.
\end{enumerate}
\end{definition}

\begin{proposition}[Obstruction theorem {\normalfont $[\ConditionalStatus]$}]
\label{prop:obstruction}
There is no universal strong ceiling theorem of the form
``$0 \leq \FC \leq 1$ for all Bianchi models under all
comparators.''  The Bianchi~IX orthogonal model under the
flat comparator provides an explicit counterexample:
$\Okaniso = -1.36 \times 10^{-3}$ from the Planck curvature
constraint gives $x < 0$, making $\FC$ undefined by
Definition~\ref{def:FC}.
\end{proposition}

\begin{proof}
By Proposition~\ref{prop:x-sign}, part~(ii), models in the
irrotational negative sector have $x < 0$.  Since
$\xmaxC > 0$ by construction, $\FC = x/\xmaxC < 0$,
violating $\FC \in [0,1]$.
\end{proof}

Table~\ref{tab:certification-classes} summarises the
certification status of each sector.

\begin{table}[t]
\caption{Certification classes for the full Bianchi
  model space (14 primary $+$ 2 growing-mode variants).
  The VER05 primary evidence table
  (Chapter~\ref{ch:results}) covers the 14 non-growing-mode
  models; the two growing-mode BVIIh variants are tested
  separately as sensitivity checks.}
\label{tab:certification-classes}
\centering\footnotesize
\setlength{\tabcolsep}{3pt}
\begin{tabular}{llccl}
\hline\hline
Sector & $\FC$ & $\FC \in [0,1]$ & $N$ & Models \\
\hline
Irrot.\ non-neg & standard & \checkmark & 9 &
  FLRW, FLRW$_{\rm t}$, BI$_{\rm o}$, BVII$_{0,\rm o}$,
  BII$_{\rm o}$, BVI$_{0,\rm o}$, BI$_{\rm t}$,
  BIII$_{\rm t}$, BVII$_{h,\rm o}$ \\
Irrot.\ negative & signed & $\times$ & 4 &
  BVIII$_{\rm o}$, BIX$_{\rm o}$, BIX$_{\rm t}$,
  BV$_{\rm t}$ \\
Vortical & blocked & $\times$ & 3 &
  BVII$_{h,\rm t}$, BVII$_{h,\rm t,g}$,
  BVII$_{h,\rm o,g}$ \\
\hline\hline
\end{tabular}
\end{table}

The certification classes are comparator-dependent: switching
from $\compflat$ to $\compmatched$ moves BVIII and BIX from
the restricted class to the certifiable class (since
$\Okaniso = 0$ by construction under $\compmatched$).
This dependence is not a weakness but an inherent feature of the
comparator-conditioned framework.


% === INLINED: figures_ch03.tex ===
% Figures for Chapter 3
% Updated: TF-D06, 2026-03-12
% Replaced: fig_MES_three_bounds (S1 corrected, hostile-audited v3),
%           fig_type_by_type_summary (Sigma2_max corrected)

% fig_4D_projection_atlas placed in ch02



% === END INLINED ===
